% 替换位置：第8.6节起至第8章结束；整体替换原8.6--8.11；8.5目标公式保留
% 这是修改片段，不是完整文档；不要整份直接追加到原文末尾。
\subsection{求解方法与最优性范围}
对8个不可拆分块进行无标签非空分区枚举，分区数分别为
\begin{equation}
\left\{\begin{matrix}8\\2\end{matrix}\right\}=127,
\qquad
\left\{\begin{matrix}8\\3\end{matrix}\right\}=966.
\end{equation}
缓存$2^8-1=255$个非空子集的资源需求和工作量后计算各分区目标。
以$12G$和$12S$进行整数比较，均衡层比较$\mathrm{CV}^2$，并列容差$10^{-12}$。
全部分区均被检查，两个基准分区在该容差下均唯一。
最优性仅针对冻结第三问基准、完整中继复制、组内可重新编号及所采用目标，
不能扩展为运输、通信与分区共同决策的全局最优性。

\subsection{分区与配置结果}
\subsubsection{两组分区}
记$\mathcal V$为全部15个服务区。两组最优分区为
\begin{align}
\mathcal V_1&=\mathcal V\setminus\{\mathrm{S008}\},&
\mathcal V_2&=\{\mathrm{S008}\},\\
\boldsymbol n_1&=(4,2,2,6,4,3,2,3),&
\boldsymbol n_2&=(0,0,1,0,0,1,1,1).
\end{align}
总需求为$\boldsymbol R^{(2)}=(4,2,3,6,4,4,3,4)$，
缺少1架C型运输机和1架中继机；中继能源组件剩余2组。
目标为$G=1$、$S=8.666667$、$\mathrm{CV}=0.735372$。

\subsubsection{三组分区}
三组最优分区为
\begin{align}
\mathcal V_1&=\mathcal V\setminus\{\mathrm{S008,S015}\},\\
\mathcal V_2&=\{\mathrm{S008}\},\qquad \mathcal V_3=\{\mathrm{S015}\}.
\end{align}
配置分别为
\begin{align}
\boldsymbol n_1&=(4,2,2,6,4,3,2,3),\\
\boldsymbol n_2&=(0,0,1,0,0,1,1,1),\\
\boldsymbol n_3&=(0,1,0,0,1,0,1,1).
\end{align}
总需求$\boldsymbol R^{(3)}=(4,3,3,6,5,4,4,5)$，
缺少1架B型运输机、1架C型运输机、1组B型电池和2架中继机；组件剩余1组。
目标为$G=2.25$、$S=10.083333$、$\mathrm{CV}=0.914171$。
\begin{table}[htbp]
\centering
\small
\caption{两种基准分区的资源需求与缺口}
\label{tab:q4_resources}
\begin{tabular}{lrrrrr}
\toprule
类型 & 库存 & 两组需求 & 两组缺口 & 三组需求 & 三组缺口 \\
\midrule
A型运输机 & 4 & 4 & 0 & 4 & 0 \\
B型运输机 & 2 & 2 & 0 & 3 & 1 \\
C型运输机 & 2 & 3 & 1 & 3 & 1 \\
A型电池 & 6 & 6 & 0 & 6 & 0 \\
B型电池 & 4 & 4 & 0 & 5 & 1 \\
C型电池 & 4 & 4 & 0 & 4 & 0 \\
中继机 & 2 & 3 & 1 & 4 & 2 \\
中继组件 & 6 & 4 & 0 & 5 & 0 \\
\bottomrule
\end{tabular}
\end{table}
\begin{table}[htbp]
\centering
\small
\caption{分区的各组工作量}
\label{tab:q4_workload}
\begin{tabular}{llrrrrrr}
\toprule
$K$ & 组 & 箱数 & 质量/kg & 体积/m$^3$ & 运输架次 & 中继任务 & 工作量/s \\
\midrule
2 & G1 & 75 & 713 & 1.882 & 23 & 3 & 55588.142 \\
2 & G2 & 5 & 45 & 0.129 & 1 & 1 & 8476.668 \\
3 & G1 & 72 & 688 & 1.815 & 22 & 3 & 53759.898 \\
3 & G2 & 5 & 45 & 0.129 & 1 & 1 & 8476.668 \\
3 & G3 & 3 & 25 & 0.067 & 1 & 1 & 8105.209 \\
\bottomrule
\end{tabular}
\end{table}

\subsection{资源缺口与额外配置的来源}
\subsubsection{组间错峰复用损失}
两组中，大组在$t=0$由P003与P005同时占用两架C型机，S008组的P012另需一架，
组间禁止调配导致总需求3架。大组由R001、R002重叠占用两架中继机，
S008组另需R001完整副本，因此中继需求也是3架。
三组再独立配置S015组，需要额外保有B型机、B型电池及中继机。
各组峰值不必同刻出现；独立执行要求求各组峰值之和，而不是分组后的全局峰值。

\subsubsection{库存剩余、额外配置与任务重复}
库存剩余为$U_\gamma=\max(0,I_\gamma-R_\gamma)$，
额外配置为$A_\gamma=R_\gamma-R^0_\gamma$，
完整中继复制引起的新增任务与能耗另行统计，三者不能混称冗余。
两组相较未分区增加1架C型机、1架中继机及1组组件；
三组增加1架B型机、1架C型机、1组B型电池、2架中继机和2组组件。
增加组件配置不等于组件库存短缺。

两组对R001额外复制一次，三组对R001额外复制两次，故中继任务总数分别为4次与5次。
复制能耗分别为
\begin{align}
E_{\mathrm{copy}}^{(2)}&=1.717471065\,\mathrm{kWh},\\
E_{\mathrm{copy}}^{(3)}&=3.434942131\,\mathrm{kWh}.
\end{align}
对应新增机体占用分别为$6276.966104\,\mathrm{s}$和$12553.932208\,\mathrm{s}$，
新增能源组件占用分别为$7174.688724\,\mathrm{s}$和$14349.377449\,\mathrm{s}$。
完整复制保持物理任务时刻，故两种分区仍有$D=0$、
$C_{\mathrm{joint}}=6302.230771\,\mathrm{s}$；运输与中继总能耗分别为
\begin{align}
E_{\mathrm{total}}^{(2)}&=74.152577548\,\mathrm{kWh},\\
E_{\mathrm{total}}^{(3)}&=75.870048613\,\mathrm{kWh}.
\end{align}

\subsection{目标优先级与解释口径比较}
\subsubsection{配置规模与工作量均衡}
将$G\to S\to\mathrm{CV}$改为$G\to\mathrm{CV}\to S$，两种$K$均选中相同基准分区。
若将均衡置于首位，则可降低工作量差异，但资源缺口明显增加。
\begin{table}[htbp]
\centering
\small
\caption{冻结方案的目标优先级比较}
\label{tab:q4_priority}
\begin{tabular}{clrrr}
\toprule
$K$ & 排序 & $G$ & $S$ & $\mathrm{CV}$ \\
\midrule
2 & $G\to S\to\mathrm{CV}$ & 1.000000 & 8.666667 & 0.735372 \\
2 & $G\to\mathrm{CV}\to S$ & 1.000000 & 8.666667 & 0.735372 \\
2 & $\mathrm{CV}\to G\to S$ & 4.500000 & 12.500000 & 0.005234 \\
3 & $G\to S\to\mathrm{CV}$ & 2.250000 & 10.083333 & 0.914171 \\
3 & $G\to\mathrm{CV}\to S$ & 2.250000 & 10.083333 & 0.914171 \\
3 & $\mathrm{CV}\to G\to S$ & 7.333333 & 15.333333 & 0.107693 \\
\bottomrule
\end{tabular}
\end{table}

\subsubsection{禁止复制中继架次}
若禁止复制，则同一中继实际保障的运输任务也必须同组。
将这些绑定加入运输关系并求传递闭包后，15个服务区合并为一个不可拆分集合。
例如R001和R002共同关联包含S001、S002、S009、S010、S012的运输块，
R003关联的S004、S005也已与R001关联；三类中继任务的绑定连通，覆盖全部服务区。
因此该解释下不存在$K=2$或$K=3$的合法非空分区，
这是禁止复制解释下的分区不可行，不是原题一般意义上的无解。

\subsubsection{继承原实体编号}
按每组出现的不同原编号计数，并对跨组同名编号分别配置实体，不能跨组共享。
本文固定基准分区的继承编号计数，与组内重新编号的最小资源向量一致；
在继承编号口径下重新优化分区，也得到相同基准需求。
这仅是本文冻结时序的结果，不说明重新编号在一般情况下没有作用。

\subsection{独立事件回放与结果有效性}
独立程序重读原始参数，复算货箱交付、能耗、SOC、充电和周转，并回放组内资源事件。
15个服务区无遗漏或重复，同架次访问点同组，80箱及其交付时刻保持不变，
16箱医疗物资和30箱首批保障物资硬时限全部满足，最小余量$194.866445\,\mathrm{s}$。
248条实际通信保障区间被继承，中继副本完整，组内资源无冲突且不跨组共享。
共12个导出场景通过检查，另以独立枚举核对127个两组分区与966个三组分区，得到10项目标核对证书。
冻结输入SHA-256在检查前后保持一致。
时间、能耗、SOC、质量、体积容差分别为$10^{-7}\,\mathrm{s}$、$10^{-8}\,\mathrm{kWh}$、
$10^{-10}$、$10^{-9}\,\mathrm{kg}$、$10^{-12}\,\mathrm{m^3}$。
通过检查表示所列配置可以承接冻结任务，不表示现有库存足够；两种$K$均无零缺口分区。
通信只继承第三问的数值验证范围，相同中继副本未额外认证共址干扰与实际网络容量。

\subsection{结果分析}
冻结方案下，两组的资源缺口、配置规模、工作量变异系数和复制能耗均低于三组，
因此按既定字典序目标优先采用两组。
分区资源代价主要来自独立保有设备后失去跨组错峰复用机会，
以及同一中继任务被多个组使用时所需的完整任务副本。
上述结论以第三问的任务时序和实际通信保障关系为固定输入，
依赖完整复制、组内同型号资源可重新编号等工作假设。
在此口径下，现有库存均不能独立承接两种分区，需要按资源类型补足缺口。
